\documentclass[a4paper,12pt]{article}
\usepackage[T2A]{fontenc}
\usepackage[utf8]{inputenc}
\usepackage[russian]{babel}
\usepackage{geometry}
\usepackage{listings}
\usepackage{longtable}
\geometry{margin=2.5cm}
\lstset{basicstyle=\ttfamily\small,breaklines=true,columns=fullflexible}

\begin{document}

\begin{titlepage}
\centering
{\large Национальный исследовательский университет ИТМО\par}
{\large Факультет программной инженерии и компьютерной техники\par}
\vspace{4cm}
{\large Отчёт к практическому заданию №2\par}
{\large по дисциплине «Системное программное обеспечение 2»\par}
\vfill
\begin{flushright}
Студент: Горинов Даниил Андреевич\\
Группа: P4116
\end{flushright}
\vfill
{Санкт-Петербург}\par
{2025/2026 учебный год\par}
\end{titlepage}

\clearpage
\pagenumbering{arabic}

\section{Цели}
Изучить механизмы синхронизации, реализовать объект последовательного потока данных и показать корректность его работы на тестовой программе по варианту.

Для варианта 4 используются SQL-выборки лабораторной работы №3, режим потоков — синхронный неблокирующий. Выборки представлены как map-reduce pipeline из независимых обработчиков, связанных последовательными byte streams.

\section{Задачи}
\begin{enumerate}
\item Изучить внутреннюю логику синхронных потоков и пассивного ожидания.
\item Модифицировать планировщик из SPO6 для поддержки ожидания готовности группы потоковых событий.
\item Описать SQL-выборки варианта и проверить их на базе PostgreSQL \texttt{ucheb}.
\item Разложить каждую выборку на этапы \texttt{source}, \texttt{filter}, \texttt{join}, \texttt{group/aggregate} и \texttt{formatter}.
\item Реализовать VM-демонстрацию конкурентного выполнения обработчиков.
\item Сравнить результаты VM-демонстрации с количеством строк, полученным SQL-запросами.
\end{enumerate}

\section{Описание работы}
Точные SQL-запросы находятся в \texttt{sql/variant\_lab3.sql}. Проверка выполнялась на базе \texttt{ucheb}; снимок результата сохранён в \texttt{sql/validation\_snapshot.md}.

VM-демонстрация находится в \texttt{tests/sql\_pipeline\_demo.asm}. В ней каждая выборка представлена набором обработчиков. Обработчики получают элементы через входные byte streams и пишут результат в выходные byte streams. Неблокирующая операция чтения или записи при отсутствии готового события возвращает состояние \texttt{would-block}; после этого обработчик переводится в \texttt{GROUP\_WAIT}, а планировщик отдаёт квант другому обработчику.

Запуск:
\begin{lstlisting}
make -C SPO7 remote-demo
\end{lstlisting}

Команда запускает \texttt{Portable.RemoteTasks.Manager.exe} через \texttt{ExecuteBinaryWithIo}, чтобы работали \texttt{SimpleClock} и \texttt{SimplePic}, затем проверяет \texttt{results/sql\_pipeline\_demo.stdout.txt} скриптом \texttt{tools/check\_sql\_pipeline\_output.py}.

Основные файлы:
\begin{itemize}
\item \texttt{tests/sql\_pipeline\_demo.asm} — модель SQL map-reduce pipeline;
\item \texttt{sql/variant\_lab3.sql} — запросы варианта;
\item \texttt{sql/validation\_snapshot.md} — проверенные количества строк;
\item \texttt{spo7.target.pdsl} и \texttt{devices.xml} — целевая VM и устройства таймера;
\item \texttt{Validation.md} — команды проверки.
\end{itemize}

\section{Аспекты реализации}
Синхронный stream не хранит большую очередь данных: передача элемента считается выполненной, когда producer и consumer согласованы. В неблокирующем режиме потоковая операция не занимает CPU в ожидании. Если вход пуст или выходной stream не готов, обработчик фиксирует \texttt{GROUP\_WAIT}, сохраняет своё состояние и возвращает управление scheduler runtime.

Планировщик использует схему из SPO6: \texttt{SimpleClock} формирует timer IRQ, \texttt{SimplePic} передаёт управление обработчику, а обработчик выполняет один шаг текущего processor stage. За счёт этого разные этапы одной выборки продвигаются независимо, а ожидание готовности stream-событий остаётся пассивным.

Декомпозиция выборок построена в терминах map-reduce:
\begin{itemize}
\item parser/source читает таблицу;
\item filter пропускает записи по условию;
\item join соединяет два входных потока по ключу;
\item group/aggregate считает \texttt{COUNT}, \texttt{AVG}, \texttt{HAVING} и устраняет дубликаты через группировку;
\item formatter пишет итоговый поток результата.
\end{itemize}

Количество обработчиков и ожидаемых \texttt{GROUP\_WAIT} задаётся таблицами в VM-программе. Поэтому степень параллелизма меняется данными сценария, а не переписыванием логики планировщика.

\section{Результаты}
Последний зафиксированный RemoteTasks-прогон:
\begin{lstlisting}
assemble = d0526229-a670-42b0-ae59-19c34a10bea0
run      = 0937d3a3-2dfd-4e25-b476-12d35df3b053
\end{lstlisting}

Результат VM-демонстрации:
\begin{lstlisting}
SPO7 SQLMR
MODE SYNC-NB GROUP-WAIT
Q1 R=1 P=5 W=1
Q2 R=0 P=5 W=1
Q3 R=5004 P=4 W=1
Q4 R=40 P=6 W=2
Q5 R=85 P=7 W=2
Q6 R=0 P=8 W=2
Q7 R=0 P=5 W=1
IRQ=48 GW=10
OK
\end{lstlisting}

Обозначения: \texttt{R} — число строк результата SQL-запроса в PostgreSQL, \texttt{P} — число обработчиков в pipeline, \texttt{W} — число событий группового ожидания. Суммарно выполнено \texttt{48} timer-driven шагов и \texttt{10} passive group waits.

\begin{center}
\begin{tabular}{|l|p{9cm}|r|}
\hline
Запрос & Смысл & Строк \\
\hline
Q1 & \texttt{RIGHT JOIN} типов ведомостей и ведомостей & 1 \\
Q2 & \texttt{INNER JOIN} людей, обучений и учеников & 0 \\
Q3 & число пар фамилия-имя после группировки & 5004 \\
Q4 & планы с двумя группами на кафедре ВТ & 40 \\
Q5 & студенты группы 4100 со средней оценкой выше средней группы 1101 & 85 \\
Q6 & студенты первого курса заочной формы до 2012-09-01 & 0 \\
Q7 & число отличников группы 3100 & 0 \\
\hline
\end{tabular}
\end{center}

\section{Выводы}
Последовательные byte streams и групповое пассивное ожидание позволяют представить SQL-выборки как набор конкурентных обработчиков. VM-демонстрация воспроизводит количества строк, полученные PostgreSQL-запросами, и показывает, что при \texttt{would-block} обработчик не ждёт активно, а уступает квант другим этапам pipeline. Это соответствует требуемому синхронному неблокирующему режиму варианта.

\end{document}
